\documentclass[11pt]{article}

% Change "review" to "final" to generate the final (sometimes called camera-ready) version.
% Change to "preprint" to generate a non-anonymous version with page numbers.
\usepackage[review]{acl}

% Standard package includes
\usepackage{times}
\usepackage{latexsym}

% For proper rendering and hyphenation of words containing Latin characters (including in bib files)
\usepackage[T1]{fontenc}
% For Vietnamese characters
% \usepackage[T5]{fontenc}
% See https://www.latex-project.org/help/documentation/encguide.pdf for other character sets

% This assumes your files are encoded as UTF8
\usepackage[utf8]{inputenc}

% This is not strictly necessary, and may be commented out,
% but it will improve the layout of the manuscript,
% and will typically save some space.
\usepackage{microtype}

% This is also not strictly necessary, and may be commented out.
% However, it will improve the aesthetics of text in
% the typewriter font.
\usepackage{inconsolata}

%Including images in your LaTeX document requires adding
%additional package(s)
\usepackage{graphicx}
\usepackage{multirow}
\usepackage{url}

% If the title and author information does not fit in the area allocated, uncomment the following
%
%\setlength\titlebox{<dim>}
%
% and set <dim> to something 5cm or larger.

\title{Modeling Bilingualism in South Asian languages}

% Author information can be set in various styles:
% For several authors from the same institution:
% \author{Author 1 \and ... \and Author n \\
%         Address line \\ ... \\ Address line}
% if the names do not fit well on one line use
%         Author 1 \\ {\bf Author 2} \\ ... \\ {\bf Author n} \\
% For authors from different institutions:
% \author{Author 1 \\ Address line \\  ... \\ Address line
%         \And  ... \And
%         Author n \\ Address line \\ ... \\ Address line}
% To start a separate ``row'' of authors use \AND, as in
% \author{Author 1 \\ Address line \\  ... \\ Address line
%         \AND
%         Author 2 \\ Address line \\ ... \\ Address line \And
%         Author 3 \\ Address line \\ ... \\ Address line}

\author{First Author \\
  Affiliation / Address line 1 \\
  Affiliation / Address line 2 \\
  Affiliation / Address line 3 \\
  \texttt{email@domain} \\\And
  Second Author \\
  Affiliation / Address line 1 \\
  Affiliation / Address line 2 \\
  Affiliation / Address line 3 \\
  \texttt{email@domain} \\}

%\author{
%  \textbf{First Author\textsuperscript{1}},
%  \textbf{Second Author\textsuperscript{1,2}},
%  \textbf{Third T. Author\textsuperscript{1}},
%  \textbf{Fourth Author\textsuperscript{1}},
%\\
%  \textbf{Fifth Author\textsuperscript{1,2}},
%  \textbf{Sixth Author\textsuperscript{1}},
%  \textbf{Seventh Author\textsuperscript{1}},
%  \textbf{Eighth Author \textsuperscript{1,2,3,4}},
%\\
%  \textbf{Ninth Author\textsuperscript{1}},
%  \textbf{Tenth Author\textsuperscript{1}},
%  \textbf{Eleventh E. Author\textsuperscript{1,2,3,4,5}},
%  \textbf{Twelfth Author\textsuperscript{1}},
%\\
%  \textbf{Thirteenth Author\textsuperscript{3}},
%  \textbf{Fourteenth F. Author\textsuperscript{2,4}},
%  \textbf{Fifteenth Author\textsuperscript{1}},
%  \textbf{Sixteenth Author\textsuperscript{1}},
%\\
%  \textbf{Seventeenth S. Author\textsuperscript{4,5}},
%  \textbf{Eighteenth Author\textsuperscript{3,4}},
%  \textbf{Nineteenth N. Author\textsuperscript{2,5}},
%  \textbf{Twentieth Author\textsuperscript{1}}
%\\
%\\
%  \textsuperscript{1}Affiliation 1,
%  \textsuperscript{2}Affiliation 2,
%  \textsuperscript{3}Affiliation 3,
%  \textsuperscript{4}Affiliation 4,
%  \textsuperscript{5}Affiliation 5
%\\
%  \small{
%    \textbf{Correspondence:} \href{mailto:email@domain}{email@domain}
%  }
%}

\begin{document}
\maketitle
\begin{abstract}
This document is a supplement to the general instructions for *ACL authors. It contains instructions for using the \LaTeX{} style files for ACL conferences.
The document itself conforms to its own specifications, and is therefore an example of what your manuscript should look like.
These instructions should be used both for papers submitted for review and for final versions of accepted papers.
\end{abstract}

\section{Introduction}

The study of language acquisition in India has drawn the attention of research from a variety of disciplines such as, psychology and education, but has been less explored from a computational perspective. Indeed, the language diversity of native speakers in India presents a complex challenge with hundreds of languages with thousands of dialects over five language families \cite{mohanty2006multilingualism}. These languages also diverse in linguistic and computational resources. While many of these languages are considered low-resource from the computational perspective. For example, Hindi has considerably larger resources compared with Telugu, though both are official state languages in certain regions of India \cite{mohanty2010languages,philip2021revisiting}. Interestingly, English appears as a constant given its use for educational settings \cite{mohanty2006multilingualism}.

Recently, there has been a growing interest in computational modeling for Indian languages, including development of Large Language Models (LLMs) optimized for these language families (e.g., IndicBERT, MuRIL) \cite{dabre2022indicbart,khanuja2021muril} and workshops dedicated to the development of models and resources for these languages CITE. While most work on small language models focused on English (see \cite{charpentier2025babylm,hu2024findings}), several studies have focused on languages other than English, including Chinese, Japanese, Solvak, and Hebrew \cite{gao2025bliss,krivs2025slovakbabylm,gelboim2025tafberta}. 

\citep{choshen2026babylm} added non-English datasets to the BabyLM challenge for Chinese and Dutch. These dataset rely on the availability of sufficient data in these two languages to create a complementary resource to the English data, i.e., 100M words spanning across data from diverse genera. We ask whether training data can be created artificially for low-resource languages directly from the English data. 

We use gpt-5 mini to translate the 100M English data provided by BabyLM workshop into Hindi and Telugu \cite{charpentier2025babylm}. As a preliminary analysis, we focus on Hindi and Telugu out of the many Indic languages. Hindi, an Indo-Aryan language  written in Devanagari, is one of the most widely spoken languages in India, while Telugu, , a Dravidian language with a distinct grammatical structure and writing system, provides a complementary typological perspective that remains significantly less resources in computational studies and resources. While we do not rely on resources to create the data, the difference in low-resources status may still effect the quality of translation, which we consider in our evaluation.

We train both monolingual and bilingual models on each language and each pairing with English. We test the models on their ability to represent the translated language, original data in each language, and developmentally-grounded measures as provided by BabyLM \cite{choshen2026babylm}. Out results show...


\section{Related Work}

Language acquisition parg



studies of Hindi and Telugu acquisition, Benchmarks of Indic LLMs and translation specifically

BLISS - multilingual baseline. M-BLIMP , use of perplexity in other studies. \cite{tampier2025recombitext} generation

\section{Methods}

We use the English 100M words data (NAME) provided by the BabyLM workshop \cite{choshen2026babylm}. We translate each of the sub-sections of the data: CHILDES \cite{macwhinney2000childes}, British National Corpus (BNC) \cite{bnc2007}, dialogue portion, Project Gutenberg (children’s stories) \cite{gerlach2020standardized}, OpenSubtitles \cite{lison-tiedemann-2016-opensubtitles2016}, Simple English Wikipedia \cite{simplewiki2023}, and Switchboard Dialog Act Corpus \cite{stolcke2000dialogue}. We translate the data using GPT-VERSION in the data to Hindi and Telugu, independently, using the following prompt: ``PROMPT''. We use BATCH AND OTHER PARAMETERIZATION. Since the training data is a translation of the original data release, the models in all languages theoretically represent equivalent semantic input that varies in other linguistic properties such as syntactic and morphological structure. We note that the resulting Hindi translation has XXX words into XXX bytes, and the Telugu translation has XXX words into XXX bytes. These translations amount to a XXX:XXX word-to-byte ratio for Hindi, and XXX:XXX for Telugu. In comparison, the Hindi and Telugu sides of the parallel English-Hindi and English-Telugu OpenSubtitles data we extracted had 1:11.97 and 1:17.18 word-to-byte ratios, respectively \citep{lison-tiedemann-2016-opensubtitles2016}.

We train monolingual and bilingual models using the code provided by the BabyLM workshop \cite{choshen2026babylm}, including GPT-WEE and GPT-BERT CITATION+NAME-FIX. The monolingual models are trained with a single language data. To analyze the effect of language settings on the resulting model, we create two versions of bilingual models for each language pair. We combine translated data from CHILDES and Gutenberg, while taking the English data for the remaining subsets. This selection aims to represent Indic native language use at home setting and English use in educational settings as observed in naturalistic settings \cite{mohanty2006multilingualism}

We train each model XXX times
for each setup with different seeds, for a total of
XXX models. The parameters for training are a batch
size of XXX, max steps of XXX with evaluation
every XXX steps, XXX warmup steps, XXX gradient
accumulation steps, and a learning rate of XXX.

We evaluate the models using two goals, analyze the models' ability to capture natural language as represented in a corpora, and their linguistic knowledge level. For language representation, we use perplexity over several sources of data, (1) the test portion parallel to the training data, (2) OpenSubtitles data from parallel curated subtitles, and (3) essays written by native-speakers of English, Hindi, and Telugu. The test data holds the highest resemblance to the training data, which provides an evaluation venue to compare the ability of small language models to learn each of the languages given their typological differences. Second, the OpenSubtitles data offers an independent source that has not been obtained through translation. Together with the first measure, it can show generalization abilities and general-language representation. 
Finally, we use the essays to evaluate the bilingual models on their ability to capture the unique properties of Indic-language native speakers use of English.

In addition, we use the measure to evaluate linguistic abilities made available by BabyLM that include benchmark for either Hindi, Telugu or both. As such, we evaluate the monolingual and bilingual models using MBLiMP (Hindi only), which measures grammatical and syntactic knowledge through minimal sentence pairs; SIB200 (Hindi and Telugu), which evaluates document classification and semantic understanding across languages; and MuBench (Hindi and Telugu), which assesses multilingual language understanding on a range of downstream NLP tasks.


\section{Results}

We evaluate the monolingual and bilingual language models using complementary intrinsic and extrinsic benchmarks designed to assess both language modeling performance and cross-lingual generalization. We first present perplexity results over all models and data sources, followed by an evaluation on linguistic abilities using the M-BLIMP benchmark.

\subsection{Perplexity}

Table~\ref{tb:mono} shows the results for the monolingual models on each of the data sources, (1) test data, (2) non-translated data, and (3) essays.



\begin{table}[t]
\centering
\footnotesize
%\setlength{\tabcolsep}{3pt}
%\renewcommand{\arraystretch}{1.0}
%\resizebox{\columnwidth}{!}{%
\begin{tabular}{lccc}
\hline
Model& Test & OS-data & Essay \\
\hline
EN-GPT-Wee & 140.50 & 306.06 & XXX \\
HI-GPT-Wee & 85.41 & 707.13 & XXX \\
TE-GPT-Wee & 134.25 & 1343.92 & XXX \\
\hline
EN-GPT-BERT & 4.8953 & XXX & XXX \\
HI-GPT-BERT & 1.3261 & XXX & XXX \\
TE-GPT-BERT & XXX & XXX & XXX \\
\hline
\end{tabular}%}
\caption{Results for the models trained on a single languages. Each column corresponds to testing on the data of the language used to train the model. GPT-BERT perplexity is a pseudo-perplexity (bidirectional, PLL-based) and is not directly comparable to the causal-LM perplexity reported for GPT-Wee.}
\label{tb:mono}
\end{table}


\begin{table*}[t]
\centering
\footnotesize
%\setlength{\tabcolsep}{3pt}
%\renewcommand{\arraystretch}{1.0}
%\resizebox{\columnwidth}{!}{%
\begin{tabular}{lccccccccc}
\hline
\multirow{2}{*}{Model} & \multicolumn{3}{c}{English} & \multicolumn{3}{c}{Hindi} &  \multicolumn{3}{c}{Telugu}\\
\cline{2-4} \cline{5-7} \cline{8-10}
 & Test & OS-data & Essay & Test & OS-data & Essay & Test & OS-data & Essay \\
\hline
EN-HI-GPT-Wee & 233.69 & 459.80 & XXX & 93.05 & 879.17 & XXX & -- & -- & -- \\
EN-TE-GPT-Wee & 199.26 & 377.26 & XXX & -- & -- & -- & 270.56 & 3937.21 & XXX \\
EN-HI-GPT-BERT & XXX & XXX & XXX & XXX & XXX & XXX & -- & -- & -- \\
EN-TE-GPT-BERT & XXX & XXX & XXX & -- & -- & -- & XXX & XXX & XXX \\
\hline
gpt5-mini & XXX & XXX & XXX & XXX & XXX & XXX & XXX & XXX & XXX \\
Sarvam-2B & 186.58 & XXX & XXX & 210.51 & XXX & XXX & 205.81 & XXX & XXX \\
\hline
\end{tabular}
\caption{Perplexity results for bilingual models and general-purpose baselines across English, Hindi, and Telugu. ``Test'' for bilingual GPT-Wee models is the filtered test set matching that model's training sources; ``--'' marks a language the model was not trained/evaluated on.}
\label{tab:bilingual-ppl}
\end{table*}

\subsection{Linguistic Knowledge}

We perform the tasks that include Hindi and/or Telugu in the basline: M-Blimp (Hindi only), SIB200 (Hindi and Telugu), and MuBench (Hindi and Telugu).

\begin{table*}[t]
\centering
\footnotesize
%\setlength{\tabcolsep}{3pt}
%\renewcommand{\arraystretch}{1.0}
%\resizebox{\columnwidth}{!}{%
\begin{tabular}{lccccccccc}
\hline
\multirow{2}{*}{Model} & \multicolumn{3}{c}{English} & \multicolumn{3}{c}{Hindi} &  \multicolumn{3}{c}{Telugu}\\
\cline{2-4} \cline{5-7} \cline{8-10}
 & MBLIMP & SIB200 & MuBench & MBLIMP & SIB200 & MuBench & MBLIMP & SIB200 & MuBench \\
\hline
EN-GPT-Wee & 0.7469 & 0.3284 & 0.2762 & XXX & XXX & XXX & XXX & XXX & XXX \\
HI-GPT-Wee & XXX & XXX & XXX & 0.9261 & 0.2353 & 0.2750 & XXX & XXX & XXX \\
TE-GPT-Wee & XXX & XXX & XXX & XXX & XXX & XXX & XXX & 0.2108 & 0.2832 \\
\hline
EN-GPT-BERT & 0.7978 & 0.2451 & XXX & XXX & XXX & XXX & XXX & XXX & XXX \\
HI-GPT-BERT & XXX & XXX & XXX & 0.9399 & 0.1961 & XXX & XXX & XXX & XXX \\
TE-GPT-BERT & XXX & XXX & XXX & XXX & XXX & XXX & XXX & XXX & XXX \\
\hline
EN-HI-GPT-Wee & 0.7070 & 0.2647 & 0.2827 & 0.9150 & 0.1324 & 0.2729 & XXX & XXX & XXX \\
EN-TE-GPT-Wee & 0.7136 & 0.2794 & 0.2803 & XXX & XXX & XXX & XXX & 0.2304 & 0.2765 \\
EN-HI-GPT-BERT & 0.7868 & 0.2500 & XXX & 0.9136 & XXX & XXX & XXX & XXX & XXX \\
EN-TE-GPT-BERT & 0.7991 & 0.2500 & XXX & XXX & XXX & XXX & XXX & XXX & XXX \\
\hline
Sarvam-2B & 0.8006 & 0.4706 & 0.2833 & 0.9772 & 0.4510 & 0.2854 & XXX & 0.3824 & 0.2768 \\
Llama 3.2 1B & 0.8188 & 0.6275 & XXX & XXX & XXX & XXX & XXX & XXX & XXX \\
\hline
\end{tabular}
\caption{Linguistic knowledge results (BLiMP/M-BLiMP, SIB-200, MuBench) across English, Hindi, and Telugu. The English MBLIMP column reports the 67-task BLiMP macro average (no multilingual M-BLiMP config exists for English); Telugu MBLIMP is left blank as no Telugu M-BLiMP benchmark exists.}
\label{tab:linguistic}
\end{table*}

\subsection{Appendices}

Use \verb|\appendix| before any appendix section to switch the section numbering over to letters. See Appendix~\ref{sec:appendix} for an example.


\section*{Limitations}

This document does not cover the content requirements for ACL or any
other specific venue.  Check the author instructions for
information on
maximum page lengths, the required ``Limitations'' section,
and so on.

\section*{Acknowledgments}

This document has been adapted
by Steven Bethard, Ryan Cotterell and Rui Yan
from the instructions for earlier ACL and NAACL proceedings, including those for
ACL 2019 by Douwe Kiela and Ivan Vuli\'{c},
NAACL 2019 by Stephanie Lukin and Alla Roskovskaya,
ACL 2018 by Shay Cohen, Kevin Gimpel, and Wei Lu,
NAACL 2018 by Margaret Mitchell and Stephanie Lukin,
Bib\TeX{} suggestions for (NA)ACL 2017/2018 from Jason Eisner,
ACL 2017 by Dan Gildea and Min-Yen Kan,
NAACL 2017 by Margaret Mitchell,
ACL 2012 by Maggie Li and Michael White,
ACL 2010 by Jing-Shin Chang and Philipp Koehn,
ACL 2008 by Johanna D. Moore, Simone Teufel, James Allan, and Sadaoki Furui,
ACL 2005 by Hwee Tou Ng and Kemal Oflazer,
ACL 2002 by Eugene Charniak and Dekang Lin,
and earlier ACL and EACL formats written by several people, including
John Chen, Henry S. Thompson and Donald Walker.
Additional elements were taken from the formatting instructions of the \emph{International Joint Conference on Artificial Intelligence} and the \emph{Conference on Computer Vision and Pattern Recognition}.


% Bibliography entries for the entire Anthology, followed by custom entries
%\bibliography{anthology,custom}
% Custom bibliography entries only
\bibliography{custom}

\appendix

\section{Example Appendix}
\label{sec:appendix}

This is an appendix.

\end{document}